\chapter{Ordered operations and a cyclic pairing}
\label{chap:cyclic-mixed-diagonal}

A bilinear form turns an operation with several inputs into a scalar
expression by pairing its output with one more input. For such expressions
to describe an action, moving the last input to the front must give the
same scalar, with the sign prescribed by the degrees. An ordered
multiplication need not have this property in its original coordinates.
Changing the coordinates of the operations, rather than discarding their
order, can reveal the required pairing.

We shall carry this out for
\[
 W(x,u)=x^3+xu^2.
\]
The mixed term forces us to retain the product of the two primitive
generators as an independent input. That input produces operations which
are absent from the list of derivatives on single generators. The resulting
four-state algebra has a nonzero operation with six inputs. Its cyclic
pairing then produces an explicit matrix Hamiltonian.

\section{The diagonal and its four states}
\label{sec:cmc-diagonal}

Let \(k\) be a field of characteristic zero. Put
\[
 R=k[[y,z]],\qquad S=R[[r,s]],\qquad x=y+r,\quad u=z+s.
\]
The letters \(x,u,y,z\) are coordinates of the potential. They are not
spacetime coordinates. All modules have parity grading.

Take the exterior algebra on two odd vectors
\(\varepsilon_1,\varepsilon_2\). Write \(\beta_i\) for
left exterior multiplication by its \(i\)-th vector, and \(\alpha_i\)
for contraction with that vector's dual. On an ordered exterior monomial,
contraction removes the \(i\)-th vector with sign \((-1)^{j-1}\) when
it occupies position \(j\), and gives zero when it is absent. Thus
\[
 \alpha_i\alpha_j+\alpha_j\alpha_i=0,\qquad
 \beta_i\beta_j+\beta_j\beta_i=0,\qquad
 \alpha_i\beta_j+\beta_j\alpha_i=\delta_{ij}.
\]
On the free \(S\)-module with basis
\(1,\varepsilon_1,\varepsilon_2,\varepsilon_1\varepsilon_2\), define the odd operator
\[
 D=r\alpha_1+s\alpha_2+
 (x^2+xy+y^2+u^2)\beta_1+y(u+z)\beta_2.
\]
The preceding relations give
\[
 D^2=(W(x,u)-W(y,z))1.
\]
Consequently \(E=\operatorname{End}_S(S\otimes\Lambda(k^2))\) is a
differential algebra with
\(\delta A=DA-(-1)^{|A|}AD\): expanding the square of this
commutator gives zero because \(D^2\) is scalar.

Differentiate while keeping \(x,u\) fixed, and set
\[
 \eta_1=-D_y=\alpha_1-(r+3y)\beta_1-(s+2z)\beta_2,\qquad
 \eta_2=-D_z=\alpha_2-y\beta_2.
\]
Their differential and multiplication are
\begin{gather}
 \delta\eta_1=3y^2+z^2,\quad \delta\eta_2=2yz,\quad
 \delta\beta_1=r,\quad\delta\beta_2=s,\label{eq:cmc-differential}\\
 \eta_1^2=-r-3y,\quad \eta_2^2=-y,\quad
 \eta_1\eta_2+\eta_2\eta_1=-s-2z,\quad
 \eta_i\beta_j+\beta_j\eta_i=\delta_{ij}.
 \label{eq:cmc-clifford}
\end{gather}
Ordered monomials \(\eta_I\beta_J\), with each index used at most
once in each increasing list, form an \(S\)-basis of \(E\).
Indeed the corresponding \(\alpha_I\beta_J\) give a basis of the
sixteen matrix operators: applying them to the four exterior basis
vectors solves for their coefficients successively. Replacing
\(\alpha_i\) by \(\eta_i\) is triangular in the number of contraction
operators, with diagonal coefficients one.

Let \(K\) be the free \(R\)-module on \(1,\xi_1,\xi_2,\xi_{12}\),
with parities \(0,1,1,0\), and differential
\[
 d1=0,\quad d\xi_1=A,\quad d\xi_2=B,\quad
 d\xi_{12}=A\xi_2-B\xi_1,\qquad
 A=3y^2+z^2,\quad B=2yz.
\]
At this point \(K\) is a complex, not an algebra with a stipulated
exterior multiplication. Define \(i:K\to E\) by the ordered
representatives \(1,\eta_1,\eta_2,\eta_1\eta_2\). Define \(p:E\to K\)
by discarding every monomial containing a \(\beta\), then setting
\(r=s=0\). Equations~\eqref{eq:cmc-differential} show that these are
chain maps and \(pi=1\).

The remaining normal coordinates have a contraction. On
\(r^a s^b y^c z^d\eta_I\beta_J\), choose the first nonzero exponent
in the ordered list \(a,b\). If neither is nonzero, set \(h=0\).
If its index is \(j\), also set \(h=0\) when \(J\) contains an
index at most \(j\). Otherwise divide by the corresponding normal
coordinate, insert \(\beta_j\) at the beginning of the \(\beta\)-list,
and multiply by \((-1)^{|I|}\). Extend by \(R\)-linearity and
coefficientwise to formal series. This gives
\begin{equation}\label{eq:cmc-contraction}
 \delta h+h\delta=1-ip,\qquad h^2=hi=ph=0.
\end{equation}
For one coordinate the formula is
\(\beta(f(r)-f(0))/r\), and its anticommutator with the differential
is \(1-ip\). For two coordinates the homotopy is
\(h_1+P_1h_2\), where \(P_1\) evaluates the first coordinate and
kills its exterior generator. Expanding its anticommutator gives
\((1-P_1)+P_1(1-P_2)=1-P_1P_2\).
The sign \((-1)^{|I|}\) makes the remaining \(\eta\)-differential
anticommute with this homotopy. The same formula proves the side
identities in~\eqref{eq:cmc-contraction}. Division lowers formal
order by at most one, so all maps are continuous.

\section{Computing the operations on composite inputs}
\label{sec:cmc-transfer}

Let \(\mathsf{s}\) denote parity-reversing suspension, and write
\[
 e=\mathsf{s}1,\quad a=\mathsf{s}\xi_1,\quad b=\mathsf{s}\xi_2,\quad c=\mathsf{s}\xi_{12}.
\]
Their parities are respectively \(1,0,0,1\).
A word means a tensor over \(R\) of these suspended states.
We take the direct sum over word lengths, completing coefficients
inside each fixed length; the \(K\)-word modules are already finite
free over \(R\). For \(E\)-words the coefficient completion includes
the normal variables in each tensor factor.
Splitting a word
at every position defines deconcatenation. An odd operation on
consecutive letters acts on a whole word with sign
\((-1)^{\text{parity of the preceding letters}}\).
These conventions determine the odd coderivation associated to
any family of operations \(b_n:(\mathsf{s}K)^{\otimes n}\to \mathsf{s}K\).

On \(\mathsf{s}E\) the differential is \(\mathsf{s}\delta\mathsf{s}^{-1}\), and the binary
operation is
\[
 B_2(\mathsf{s}U,\mathsf{s}V)=(-1)^{|U|}\mathsf{s}(UV).
\]
The contraction determines maps \(F_n:(\mathsf{s}K)^{\otimes n}\to \mathsf{s}E\)
of parity zero and operations \(b_n\) of parity one:
\begin{equation}\label{eq:cmc-recursion}
 F_1=\mathsf{s}i\mathsf{s}^{-1},\qquad
 R_n=\sum_{j=1}^{n-1}B_2(F_j\otimes F_{n-j}),\qquad
 F_n=-\mathsf{s}h\mathsf{s}^{-1}R_n,\quad b_n=\mathsf{s}p\mathsf{s}^{-1}R_n\quad(n\ge2),
\end{equation}
with \(b_1=\mathsf{s}d\mathsf{s}^{-1}\).

Here the nonzero differential on \(K\) is retained. To see directly
why the recurrence gives a morphism, extend \(F\) to words by
consecutive blocks and extend \(b\) by the insertion signs above.
Put \(d_E=\mathsf{s}\delta\mathsf{s}^{-1}\), \(I=\mathsf{s}i\mathsf{s}^{-1}\),
\(P=\mathsf{s}p\mathsf{s}^{-1}\), \(H=\mathsf{s}h\mathsf{s}^{-1}\), and let \(\mathcal B\) be the coderivation
of the binary product. If \(\pi\) projects to words of length one,
set \(\mathcal R=\pi\mathcal B F\). Then
\[
 \pi F=I\pi-H\mathcal R,\qquad
 \pi b=b_1\pi+P\mathcal R.
\]
Suppose the defect \(\mathcal D=(d_E+\mathcal B)F-Fb\) is zero on shorter
words. Its coproduct makes it a relative coderivation, so its value
on the next word has length one. Hence \(\pi\mathcal B\mathcal D=0\).
Since \(d_E\mathcal B+\mathcal B d_E=\mathcal B^2=0\), this says
\(d_E\mathcal R+\mathcal Rb=0\) there.
Substitution of the two displayed recurrences
and~\eqref{eq:cmc-contraction} now gives
\(\pi\mathcal D=H(d_E\mathcal R+\mathcal Rb)=0\).
Induction proves the morphism equation. The highest-length part
of \(F\) is \(I^{\otimes n}\), which is split injective.
Applying the target differential twice therefore gives \(b^2=0\).
The first component is the chain homotopy equivalence \(i\).
Unit inputs give their binary unit products; in longer words
the side identities \(HI=PH=H^2=0\) make their contributions zero.
Thus both the algebra and the comparison are strictly unital.

Write \(\phi_n=\mathsf{s}^{-1}F_n\), so the entries of the following table
are actual operators in \(E\). These are all its nonzero components
on nonunit words of length at least two:

\begin{longtable}{c l l}
\toprule
Arity & Input word & \(\phi_n\) \\
\midrule
\endhead
2 & \(aa\) & \(-\beta_1\) \\
2 & \(ac\) & \(\eta_2\beta_1\) \\
2 & \(ba\) & \(-\beta_2\) \\
2 & \(bc\) & \(\eta_2\beta_2\) \\
2 & \(ca\) & \(\eta_2\beta_1-\eta_1\beta_2\) \\
2 & \(cc\) & \(y\beta_1+\eta_1\eta_2\beta_2\) \\
3 & \(aca\) & \(\beta_1\beta_2\) \\
3 & \(acc\) & \(\eta_2\beta_1\beta_2\) \\
3 & \(bcc\) & \(-\beta_2\) \\
3 & \(caa\) & \(\beta_1\beta_2\) \\
3 & \(cac\) & \(\beta_1+\eta_2\beta_1\beta_2\) \\
3 & \(ccc\) & \(\eta_2\beta_1-\eta_1\beta_2\) \\
4 & \(accc\) & \(\beta_1\beta_2\) \\
4 & \(cacc\) & \(\beta_1\beta_2\) \\
4 & \(ccca\) & \(\beta_1\beta_2\) \\
4 & \(cccc\) & \(\eta_2\beta_1\beta_2+\beta_1\) \\
5 & \(ccccc\) & \(\beta_1\beta_2\) \\
\bottomrule
\end{longtable}

The table follows from~\eqref{eq:cmc-recursion} using only
\eqref{eq:cmc-clifford} and the displayed rule for \(h\).
For example, with \(C=\eta_1\eta_2\),
\[
 C^2=-y(r+3y)-(s+2z)C,\qquad
 \phi_2(c,c)=-h(C^2)=y\beta_1+C\beta_2.
\]
Every higher \(\phi_n\) lies in the left ideal
\(E\beta_1+E\beta_2\), which \(p\) kills.
Consequently, when projecting the recurrence for an operation,
only its last split can contribute:
\begin{equation}\label{eq:cmc-last-split}
 \mathsf{s}^{-1}b_n(v_1,\ldots,v_n)=
 (-1)^{1+\sum_{j<n}|v_j|}
 p\bigl(\phi_{n-1}(v_1,\ldots,v_{n-1})i(\mathsf{s}^{-1}v_n)\bigr).
\end{equation}
This makes the passage from the operator table to every ordered
operation explicit. In particular
\[
 b_6(c,c,c,c,c,c)
 =\mathsf{s}\,p(\beta_1\beta_2\eta_1\eta_2)=-e.
\]
Moving the two \(\beta\)'s past the two \(\eta\)'s gives scalar
term \(-1\); every other term still has a \(\beta\) on the right.

There is a finite bound for these particular inputs. Give \(r,s,y,z\)
weight one, each \(\eta_i\) weight \(1/2\), and each \(\beta_i\)
weight \(-1/2\). Multiplication preserves weight and \(h\) lowers
it by \(3/2\). A comparison tree with \(n\) inputs contains \(n-1\)
homotopies. Its input basis states have total weight at most \(n\),
so its output has weight at most \((3-n)/2\). Matrix basis
monomials have weight at least \(-1\). Thus \(F_n=0\) for \(n\ge6\).
An operation tree has \(n-2\) homotopies and its projected output
has nonnegative weight, giving \(b_n=0\) for \(n\ge7\).
Coefficient-linearity extends these conclusions from basis states
to every input. The tables and bounds therefore describe the entire
formal model.

\section{Choosing coordinates compatible with a pairing}
\label{sec:cmc-cyclic}

For a bilinear form \(\omega\) on \(\mathsf{s}K\), cyclicity means
\begin{equation}\label{eq:cmc-cyclicity}
 \omega\bigl(b_n(v_1,\ldots,v_n),v_{n+1}\bigr)
 =
 (-1)^{|v_1|\sum_{j=2}^{n+1}|v_j|}
 \omega\bigl(b_n(v_2,\ldots,v_{n+1}),v_1\bigr).
\end{equation}
We seek a nondegenerate even form: it pairs equal parities, and
the map to the \(R\)-linear dual that it defines is invertible.
On the critical fibre, the natural candidate has
\begin{equation}\label{eq:cmc-form}
 \omega(e,c)=\omega(c,e)=1,\qquad
 \omega(a,b)=-1,\qquad\omega(b,a)=1,
\end{equation}
and all other entries zero. This is graded antisymmetric.
The ordered transfer is not yet cyclic for it.
For instance \(b_3(b,a,b)=e\), whereas \(b_3(a,b,c)=0\);
pairing with \(c\) and \(b\) respectively violates
\eqref{eq:cmc-cyclicity}.

Define a coalgebra change of coordinates \(G\) by
\[
 G_1=1,\qquad G_2(b,c)=-e,\qquad G_2(c,b)=e,
\]
with every other higher component zero. Its inverse has first
component one and second component \(-G_2\).
Indeed a nested \(G_2\) encounters its preceding unit output and
vanishes. Let \(b^G=G^{-1}bG\). This preserves the differential,
the unit, and the full homotopy type.

Over \(R\), the form for this intermediate model is
\[
 \omega_R=\omega,\quad\text{except that}\quad
 \omega_R(c,c)=-2z.
\]
The parameter term matters. For example the constant form before
this correction gives cyclicity defects \(2z,-2z,2z,2z,-2z,-2z\)
on the words \(ecc,abc,bac,cab,cba,cce\).
Set
\[
 L(c)=c+ze,\qquad L(e)=e,\quad L(a)=a,\quad L(b)=b.
\]
Then \(L^*\omega_R=\omega\): its only new diagonal value is
\(\omega_R(c+ze,c+ze)=-2z+2z=0\).
The constant-form operations are
\[
 \widehat b=L^{-1}b^G L^{\otimes\bullet}.
\]
Their complete nonunit table is as follows. In addition,
\(\widehat b_2(e,v)=v\) and
\(\widehat b_2(v,e)=(-1)^{|v|+1}v\);
all other unit-containing operations vanish.

\begin{longtable}{c l l}
\toprule
Arity & Input word & \(\widehat b_n\) \\
\midrule
\endhead
1 & \(a\) & \((3 y^{2} + z^{2})e\) \\
1 & \(b\) & \((2 y z)e\) \\
1 & \(c\) & \((3 y^{2} + z^{2})b+(- 2 y z)a\) \\
2 & \(aa\) & \((3 y)e\) \\
2 & \(ab\) & \(-c+(z)e\) \\
2 & \(ac\) & \((3 y)b+(- z)a\) \\
2 & \(ba\) & \(c+(z)e\) \\
2 & \(bb\) & \((y)e\) \\
2 & \(bc\) & \((- y)a+(z)b\) \\
2 & \(ca\) & \((3 y)b+(- z)a\) \\
2 & \(cb\) & \((- y)a+(z)b\) \\
2 & \(cc\) & \((3 y^{2} + 3 z^{2})e\) \\
3 & \(aaa\) & \(e\) \\
3 & \(aac\) & \(b\) \\
3 & \(abb\) & \(e\) \\
3 & \(abc\) & \(a\) \\
3 & \(aca\) & \(b\) \\
3 & \(acb\) & \(-a\) \\
3 & \(acc\) & \((2 y)e\) \\
3 & \(bab\) & \(-e\) \\
3 & \(bac\) & \(-a\) \\
3 & \(bba\) & \(e\) \\
3 & \(bbc\) & \(b\) \\
3 & \(bca\) & \(-a\) \\
3 & \(bcb\) & \(-b\) \\
3 & \(bcc\) & \((2 z)e\) \\
3 & \(caa\) & \(b\) \\
3 & \(cab\) & \(-a\) \\
3 & \(cac\) & \((2 y)e\) \\
3 & \(cba\) & \(a\) \\
3 & \(cbb\) & \(b\) \\
3 & \(cbc\) & \((2 z)e\) \\
3 & \(cca\) & \((2 y)e\) \\
3 & \(ccb\) & \((2 z)e\) \\
3 & \(ccc\) & \((2 y)b+(- 2 z)a\) \\
4 & \(aacc\) & \(e\) \\
4 & \(accc\) & \(b\) \\
4 & \(bbcc\) & \(e\) \\
4 & \(bccc\) & \(-a\) \\
4 & \(caac\) & \(e\) \\
4 & \(cbbc\) & \(e\) \\
4 & \(ccaa\) & \(e\) \\
4 & \(ccbb\) & \(e\) \\
4 & \(ccca\) & \(b\) \\
4 & \(cccb\) & \(-a\) \\
6 & \(cccccc\) & \(-e\) \\
\bottomrule
\end{longtable}

\begin{theorem}\label{thm:cmc-cyclic-comparison}
The table defines a strictly unital cyclic \(A_\infty\)-algebra
over \(R\), with the nondegenerate form~\eqref{eq:cmc-form}.
The comparison to the diagonal differential algebra \(E\) has
components
\[
 F^{\mathrm{cyc}}_1=F_1L,\qquad
 F^{\mathrm{cyc}}_2=F_2+F_1G_2,\qquad
 F^{\mathrm{cyc}}_n=F_n\quad(n\ge3).
\]
It is a strictly unital, continuous \(R\)-linear
\(A_\infty\)-quasi-isomorphism. The same formulas survive
specialization of the retained parameters.
\end{theorem}

\begin{proof}
For completeness, the conjugation can be computed without
inverting an infinite series of unknown maps. On a nonunit word
\(v_1\cdots v_n\), expand \(G\) into consecutive blocks of length
one or two, using its two specified binary values. If \(U_n\)
is the length-one component of \(bG\), the equation \(bG=Gb^G\)
gives, for \(n\ge2\),
\[
 b^G_n=U_n
 -G_2\bigl(b^G_{n-1}(v_1,\ldots,v_{n-1}),v_n\bigr)
 -(-1)^{|v_1|}
 G_2\bigl(v_1,b^G_{n-1}(v_2,\ldots,v_n)\bigr).
\]
Here \(b^G_1=b_1\). This recursion, \eqref{eq:cmc-last-split},
and the operator table give exactly the displayed coefficients.
Applying \(L\) adds unit inputs only; their contributions are
the stated binary products. The same weights used above apply
to \(G_2\), which lowers weight by \(3/2\), and to \(L\), which
preserves weight. Thus no operations beyond the table are omitted.

Cyclicity is a finite coefficient identity on this table. Pair
each displayed output with the four basis vectors using
\eqref{eq:cmc-form} and rotate the resulting word.
The paired coefficients agree with the sign in
\eqref{eq:cmc-cyclicity}; words outside these rotation orbits
give zero on both sides. For example \(bab c\) has coefficient
\(-1\), as does its rotation \(abc b\); the seven-letter word
\(c^7\) has coefficient \(-1\) and positive rotation sign.
The unary entries also obey the identity, so the nonzero
parameter differential is included. The matrix of the form
has determinant \(-1\).

The map is the composite \(FGL\).
Every \(G_2\) output is a unit, so a component \(F_j\) with
\(j\ge2\) kills it. Similarly the additional \(ze\) in \(L(c)\)
can contribute only to \(F_1\). These facts give the displayed
three formulas. Its linear component is \(iL\), a chain
homotopy equivalence. All identities are between finite
compositions in each arity and are coefficient-linear.
The explicit contraction remains an identity after
specialization; no flatness assertion about arbitrary
quasi-isomorphisms is needed.
\end{proof}

The corrected top-state representative is
\(\eta_1\eta_2+z1\). On the diagonal \(r=s=0\),
\eqref{eq:cmc-clifford} identifies it with
\((\eta_1\eta_2-\eta_2\eta_1)/2\).
Thus the constant pairing selects the skew-symmetric top
representative along the diagonal, while the higher components
retain the ordered information away from it.

The individual sixth coefficient is not determined by the
underlying \(A_\infty\)-homotopy type, even among models cyclic
for the same form. On the critical fibre, introduce a scalar
\(q\) and replace the binary component of \(G\) by
\[
 G_2(b,c)=-e,\quad G_2(c,b)=e,\quad
 G_2(a,c)=qe,\quad G_2(c,a)=-qe,\quad G_2(c,c)=qb.
\]
Every other higher component is zero and \(G_1=1\).
This remains an invertible coalgebra map by word-length
triangularity. The recursion in the proof gives cyclic
operations for the same constant form and
\[
 (G^{-1}bG)_6(c,c,c,c,c,c)=(7q^2-1)e.
\]
Both statements are polynomial coefficient identities: substitute
these five binary values in that recursion and pair each resulting
output as in~\eqref{eq:cmc-cyclicity}. Its weight argument still
excludes arities greater than six. Over a field containing
\(\sqrt7\), taking \(q=1/\sqrt7\) kills this coefficient while
preserving the full \(A_\infty\)-homotopy type. The remaining operations
change as well; the existence of the original sixth term alone
cannot prove that it is indispensable in every coordinate system.

The cyclic pairing turns the ordered operations into scalar functions of
one additional input. Their matrix action will be constructed in
Chapter~\ref{chap:cmc-matrix-action}, after the bracket and cotangent
action have been introduced.
